\begin{frame}{이번 주차 정리}
  \begin{itemize}
    \item \kb{11.1 신뢰 영역과 확률비} —
      \begin{itemize}
        \item 신뢰 영역 = 한 걸음의 크기를 ``믿을 수 있는
          범위''(KL 발산)로 제한하는 직관.
        \item 확률비 $r_t=\pi_\theta/\pi_{\mathrm{old}}$ : 첫
          스텝에서는 정직한 policy gradient, $r_t$ 가 크면
          분산 $e^{\mu^2}-1$ 으로 지수적 불어나기 → 소수 극단
          샘플 지배.
      \end{itemize}
    \item \kb{11.2 PPO 클리핑과 연속 제어} —
      \begin{itemize}
        \item 클리핑 한 줄: min의 모서리가 \textbf{개선
          방향에만} 상한(비대칭). 전체 목적함수 = 정책 +
          Critic $-$ 엔트로피.
        \item Pendulum 실습: 연속 제어는 탐색 체적 때문에
          본질적으로 느림(12만 스텝 ≈ 8\% 개선).
      \end{itemize}
    \item \kb{11.3 GAE와 실전 스택} —
      \begin{itemize}
        \item GAE = 여러 시점 TD 오차의 지수 가중합
          ($\lambda$ 로 TD$\leftrightarrow$MC 절충, 기본값
          0.95) + 엔트로피 보너스가 탐험 소실을 늦춤.
        \item TRPO의 300$\sim$400줄 장치를 클리핑 100줄로
          대체 → ``보장은 근사로, 실용성은 그대로'' —
          표준이 된 이유.
      \end{itemize}
  \end{itemize}
  \vskip0.6em
  \begin{block}{PPO 공식}
    \kb{PPO = (11.1 확률비) $\times$ (11.2 클리핑) + (11.3 GAE·
    엔트로피·실전 스택)} — 각각 빠지면 ``폭주(REINFORCE)'',
    ``불안정(클립 없는 IS)'', ``느리고 탐험 소실(GAE/엔트로피
    없음)''.
  \end{block}
  \vskip0.3em
  \begin{block}{다음 주차 — Week 12. 모방학습과 인간 피드백}
    보상의 시행착오 대신 \textbf{사람의 시연과 선호}로부터 배우는
    방향 — Week 2$\sim$11의 보상 최적화 알고리즘 중 가장
    강한 PPO가 ``왜 보상 최적화만으론 안 되는가''를 비교하기
    위한 기준점이 됨.
  \end{block}
\end{frame}
